\textbf{Phần 1. Trắc nghiệm lựa chọn}
\setcounterpageref{ex}{0}

\begin{ex} %[3H3N21]
Khối lập phương có bao nhiêu mặt?
\choice
{$4$ mặt}
{$8$ mặt}
{\True $6$ mặt}
{$12$ mặt}
\loigiai{
Khối lập phương có 6 mặt.\\
\textbf{Đáp án của câu hỏi là C}
}
\end{ex}

\begin{ex} %[3H3N21]
Khối hộp chữ nhật có bao nhiêu đỉnh?
\choice
{$6$ đỉnh}
{\True $8$ đỉnh}
{$10$ đỉnh}
{$12$ đỉnh}
\loigiai{
Khối hộp chữ nhật có 8 đỉnh.\\
\textbf{Đáp án của câu hỏi là B}
}
\end{ex}

\begin{ex} %[3H3H21]
Khối lập phương có bao nhiêu cạnh?
\choice
{$8$ cạnh}
{$6$ cạnh}
{$10$ cạnh}
{\True $12$ cạnh}
\loigiai{
Khối lập phương có 12 cạnh bằng nhau.\\
\textbf{Đáp án của câu hỏi là D}
}
\end{ex}

\begin{ex} %[3H3H21]
Các mặt của khối lập phương đều là hình gì?
\choice
{\True Hình vuông}
{Hình chữ nhật}
{Hình tam giác}
{Hình tròn}
\loigiai{
Tất cả 6 mặt của khối lập phương đều là hình vuông bằng nhau.\\
\textbf{Đáp án của câu hỏi là A}
}
\end{ex}

\begin{ex} %[3H3H21]
Các mặt của khối hộp chữ nhật thường là hình gì?
\choice
{Hình tam giác}
{Hình tròn}
{\True Hình chữ nhật}
{Hình thoi}
\loigiai{
Các mặt của khối hộp chữ nhật là hình chữ nhật.\\
\textbf{Đáp án của câu hỏi là C}
}
\end{ex}

\begin{ex} %[3H3H21]
Một con xúc xắc chuẩn có dạng khối gì?
\choice
{\True Khối lập phương}
{Khối hộp chữ nhật}
{Khối trụ}
{Khối cầu}
\loigiai{
Con xúc xắc có dạng khối lập phương.\\
\textbf{Đáp án của câu hỏi là A}
}
\end{ex}

\begin{ex} %[3H3V21]
Một hộp quà có dạng khối lập phương có cạnh dài $5\text{ cm}$. Tổng độ dài tất cả 12 cạnh của hộp quà đó là
\choice
{$50\text{ cm}$}
{$70\text{ cm}$}
{$40\text{ cm}$}
{\True $60\text{ cm}$}
\loigiai{
$5 \times 12 = 60\text{ cm}$.\\
\textbf{Đáp án của câu hỏi là D}
}
\end{ex}

\begin{ex} %[3H3V21]
Một khối rubik $3 \times 3$ có dạng hình lập phương. Số mặt có màu khác nhau của khối rubik là
\choice
{$4$ mặt}
{$8$ mặt}
{\True $6$ mặt}
{$12$ mặt}
\loigiai{
Khối rubik có 6 mặt.\\
\textbf{Đáp án của câu hỏi là C}
}
\end{ex}

\begin{ex} %[3H3V21]
Người ta dán giấy màu vào tất cả các mặt của một hộp phấn hình lập phương. Hỏi người ta phải dán bao nhiêu mặt?
\choice
{$4$ mặt}
{\True $6$ mặt}
{$8$ mặt}
{$12$ mặt}
\loigiai{
Dán đủ 6 mặt.\\
\textbf{Đáp án của câu hỏi là B}
}
\end{ex}

\begin{ex} %[3H3V21]
Xếp 8 khối lập phương nhỏ cạnh $1\text{ cm}$ lại thành một khối lập phương lớn. Khối lập phương lớn có cạnh dài là
\choice
{\True $2\text{ cm}$}
{$1\text{ cm}$}
{$3\text{ cm}$}
{$4\text{ cm}$}
\loigiai{
Khối lập phương lớn có cạnh $2\text{ cm}$ (vì $2 \times 2 \times 2 = 8$).\\
\textbf{Đáp án của câu hỏi là A}
}
\end{ex}

\begin{ex} %[3H3V21]
Một viên gạch xây dựng thường có dạng khối gì?
\choice
{Khối lập phương}
{Khối cầu}
{Khối nón}
{\True Khối hộp chữ nhật}
\loigiai{
Viên gạch có dạng khối hộp chữ nhật.\\
\textbf{Đáp án của câu hỏi là D}
}
\end{ex}

\begin{ex} %[3H3C21]
Một khung dây thép hình lập phương có cạnh $10\text{ cm}$. Tính tổng chiều dài dây thép cần dùng để làm khung đó (bỏ qua mối nối).
\choice
{$100\text{ cm}$}
{\True $120\text{ cm}$}
{$60\text{ cm}$}
{$80\text{ cm}$}
\loigiai{
Khối lập phương có 12 cạnh: $10 \times 12 = 120\text{ cm}$.\\
\textbf{Đáp án của câu hỏi là B}
}
\end{ex}

\textbf{Phần 2. Trắc nghiệm đúng sai}
\setcounterpageref{ex}{0}

\begin{ex} %[3H3H21]
Xét tính đúng sai về khối lập phương và khối hộp chữ nhật:
\choiceTF
{\True Khối lập phương và khối hộp chữ nhật đều có 6 mặt, 8 đỉnh và 12 cạnh.}
{\True Tất cả 6 mặt của khối lập phương là các hình vuông bằng nhau.}
{Khối lập phương có 8 mặt và 6 đỉnh.}
{\True Khối lập phương có 6 mặt và 8 đỉnh.}
\loigiai{
\begin{itemize}
    \item 6 mặt, 8 đỉnh, 12 cạnh: Đúng.
    \item 6 mặt hình vuông bằng nhau: Đúng.
    \item Có 6 mặt và 8 đỉnh (không phải 8 mặt 6 đỉnh). Ý 3 Sai, Ý 4 Đúng.
\end{itemize}
\begin{itemize}
    \item a) \textbf{Đúng}.
    \item b) \textbf{Đúng}.
    \item c) \textbf{Sai}.
    \item d) \textbf{Đúng}.
\end{itemize}
}
\end{ex}

\begin{ex} %[3H3V21]
Một thùng các-tông có dạng khối hộp chữ nhật. Xét tính đúng sai:
\choiceTF
{\True Thùng có 8 đỉnh ở các góc.}
{\True Thùng có 12 cạnh.}
{Thùng có 4 mặt xung quanh và không có mặt đáy.}
{\True Thùng có tất cả 6 mặt (gồm 4 mặt bên và 2 mặt đáy).}
\loigiai{
\begin{itemize}
    \item 8 đỉnh: Đúng.
    \item 12 cạnh: Đúng.
    \item Có đủ 6 mặt (4 mặt bên và 2 đáy). Ý 3 Sai, Ý 4 Đúng.
\end{itemize}
\begin{itemize}
    \item a) \textbf{Đúng}.
    \item b) \textbf{Đúng}.
    \item c) \textbf{Sai}.
    \item d) \textbf{Đúng}.
\end{itemize}
}
\end{ex}

\textbf{Phần 3. Trả lời ngắn}
\setcounterpageref{ex}{0}

\begin{ex} %[3H3H21]
Khối lập phương có bao nhiêu đỉnh?
\shortans{8}
\loigiai{
\textbf{Đáp số:} $8$\\
8 đỉnh.
}
\end{ex}

\begin{ex} %[3H3H21]
Khối hộp chữ nhật có bao nhiêu cạnh?
\shortans{12}
\loigiai{
\textbf{Đáp số:} $12$\\
12 cạnh.
}
\end{ex}

\begin{ex} %[3H3V21]
Khối lập phương có bao nhiêu mặt?
\shortans{6}
\loigiai{
\textbf{Đáp số:} $6$\\
6 mặt.
}
\end{ex}

\begin{ex} %[3H3V21]
Một khung dây hình lập phương có cạnh dài $6\text{ cm}$. Tính tổng độ dài 12 cạnh của khung dây đó theo xăng-ti-mét.
\shortans{72}
\loigiai{
\textbf{Đáp số:} $72$\\
$6 \times 12 = 72\text{ cm}$.
}
\end{ex}

\begin{ex} %[3H3V21]
Xếp 2 khối lập phương nhỏ cạnh $3\text{ cm}$ khít cạnh nhau để tạo thành một khối hộp chữ nhật. Tính chiều dài của khối hộp chữ nhật đó theo xăng-ti-mét.
\shortans{6}
\loigiai{
\textbf{Đáp số:} $6$\\
$3 + 3 = 6\text{ cm}$.
}
\end{ex}

\begin{ex} %[3H3C21]
Người ta xếp các khối lập phương nhỏ cạnh $1\text{ cm}$ thành một khối lập phương lớn cạnh $3\text{ cm}$. Cần dùng tất cả bao nhiêu khối lập phương nhỏ?
\shortans{27}
\loigiai{
\textbf{Đáp số:} $27$\\
$3 \times 3 \times 3 = 27$ khối lập phương nhỏ.
}
\end{ex}
